\chapter{Literature Survey}

This chapter reviews the work this project builds on. It covers host-based intrusion
detection from system calls, security of containers and Kubernetes, eBPF-based
monitoring and detection, machine-learning techniques for intrusion detection, and
the datasets available for evaluation. It ends with a comparison of existing
approaches and the research gap addressed by this project.

\section{Host-Based Intrusion Detection from System Calls}

Forrest et al.~\cite{forrest1996} showed that short sequences of system calls form a
stable ``sense of self'' for Unix processes and that deviations from this normal
profile reveal intrusions. This work established system calls as the primary data
source for host-based intrusion detection systems (HIDS). Later research replaced
fixed sequence databases with statistical and machine-learning models, but the core
idea of judging a process by its interaction with the kernel remains the foundation of
runtime security today.

Creech and Hu~\cite{creech2013} created the ADFA-LD dataset of Linux system-call traces
to replace the outdated KDD data. Their experiments showed that models based only on
system-call numbers struggle to separate modern attacks from normal behaviour, which
motivates the use of richer context such as arguments, file paths and process lineage,
as done in this project.

\section{Security of Containers and Kubernetes}

Sultan et al.~\cite{sultan2019} surveyed container security and grouped threats into
attacks from inside a container, between containers, on the host and through the
image supply chain. They emphasise that containers share the host kernel, so a
compromised container is a direct threat to the node. The MITRE ATT\&CK Containers
matrix~\cite{mitre} catalogues concrete techniques such as escape to host (T1611),
deploying containers, accessing the container API and resource hijacking for
crypto-mining (T1496). These techniques were used to define the attack categories of
our dataset and to label alerts.

Abed et al.~\cite{abed2015} applied the ``bag of system calls'' representation to Linux
containers and showed that anomalous behaviour of containerised applications can be
detected from system-call frequencies. Tien et al.~\cite{tien2019} proposed KubAnomaly,
a neural-network based anomaly detector for Docker and Kubernetes that monitors
containers through a kernel-level agent. Karn et al.~\cite{karn2021} detected
crypto-mining pods in container clouds from system calls using explainable machine
learning, demonstrating that model explanations help operators trust detections.

\section{eBPF-Based Monitoring and Detection}

eBPF allows verified programs to run in the Linux kernel and attach to tracepoints,
kprobes and network hooks with low overhead~\cite{gregg2019, ebpfio}. It has become the
basis of several CNCF runtime-security tools. \textbf{Falco}~\cite{falco} evaluates a
rule language over system-call events; \textbf{Tetragon}~\cite{tetragon} provides
eBPF-based process, file and network observability with Kubernetes identity and can
enforce policies in the kernel; \textbf{Tracee}~\cite{tracee} offers signature-based
detection on top of eBPF tracing. These tools are rule-driven: they are precise for
known behaviour but do not score previously unseen variants.

Research has started to combine eBPF with machine learning. Bachl et
al.~\cite{bachl2021} implemented a flow-based IDS whose decision tree runs directly
inside eBPF. Zhang et al.~\cite{zhang2024} executed a neural network in the kernel with
eBPF for real-time intrusion detection and prevention. Kostopoulos et
al.~\cite{kostopoulos2025} presented eIDPS, an eBPF and machine-learning based network
intrusion detection and prevention solution. These systems focus mainly on
\emph{network} traffic; host-level process behaviour inside Kubernetes pods, with Pod
and workload context, is less covered.

\section{Machine Learning for Intrusion Detection}

Random Forests~\cite{breiman2001} combine many decision trees trained on bootstrap
samples and random feature subsets. They handle mixed binary and numeric features,
need little preprocessing, give calibrated class probabilities and provide feature
importances, which makes them a common and strong baseline for intrusion detection on
tabular features. Isolation Forest~\cite{liu2008} is an unsupervised anomaly detector
that isolates rare points with random splits; it needs no attack labels but cannot tell
\emph{which} attack occurred and usually produces more false positives. The first
prototype of this project used Isolation Forest, and we keep it as a baseline.

Sommer and Paxson~\cite{sommer2010} warned that machine learning for intrusion
detection often looks excellent in the laboratory but fails in operation because of the
gap between training data and real traffic, the high cost of false positives and the
lack of interpretability. We address these concerns by (i) combining ML with
interpretable rules, (ii) reporting false-positive rates explicitly, and (iii)
evaluating on real public data in addition to our synthetic data.

\section{Datasets for Runtime Intrusion Detection}

\begin{itemize}
    \item \textbf{ADFA-LD}~\cite{creech2013}: Linux system-call traces of normal use and
          attacks; only system-call numbers, no arguments or containers.
    \item \textbf{LID-DS}~\cite{grimmer2019}: system-call recordings of containerised
          applications with CVE-based attack scenarios, recorded with sysdig.
    \item \textbf{BETH}~\cite{highnam2021}: about eight million kernel events collected
          with eBPF on cloud honeypots, labelled as benign, suspicious or evil.
    \item \textbf{Splunk Attack Data / Cisco Isovalent}~\cite{splunkattack}: Tetragon
          (eBPF) process events from attack simulations in a Kubernetes cluster, mapped to
          MITRE ATT\&CK.
\end{itemize}
None of these datasets has the exact schema used by our sensor (one row per kernel
event with Kubernetes namespace, behaviour-window counters and attack category), so we
generated our own labelled dataset and converted BETH and the Isovalent data to our
schema for independent testing (Chapter~4).

\section{Comparison of Existing Approaches}

\begin{table}[H]
\centering
\caption{Comparison of related approaches}
\label{tab:litcompare}
\small
\begin{tabularx}{\textwidth}{>{\raggedright\arraybackslash}p{2.6cm} >{\raggedright\arraybackslash}p{1.9cm} >{\raggedright\arraybackslash}p{2.1cm} >{\raggedright\arraybackslash}X}
\toprule
\textbf{Work / tool} & \textbf{Data source} & \textbf{Method} & \textbf{Limitation for our goal} \\
\midrule
Forrest et al.~\cite{forrest1996} & System calls & Sequence profile & No containers, no context \\
Abed et al.~\cite{abed2015} & System calls & Bag of syscalls & Offline, no Kubernetes metadata \\
KubAnomaly~\cite{tien2019} & Kernel agent & Neural network & Unsupervised score, no attack type \\
Karn et al.~\cite{karn2021} & System calls & Explainable ML & Only crypto-mining \\
Falco / Tracee~\cite{falco,tracee} & eBPF / syscalls & Rules & Detects only what rules describe \\
Tetragon~\cite{tetragon} & eBPF & Policies & Observability / enforcement, no ML scoring \\
Bachl et al., Zhang et al.~\cite{bachl2021,zhang2024} & eBPF (network) & ML in kernel & Network flows, not process behaviour \\
eIDPS~\cite{kostopoulos2025} & eBPF (network) & ML + prevention & Network IDS, not pod runtime \\
\midrule
\textbf{This project} & eBPF (16 syscalls) & RF + 19 rules & Single-node evaluation; synthetic training attacks \\
\bottomrule
\end{tabularx}
\end{table}

\section{Research Gap}

From the survey we identified the following gap: there is no simple, open,
end-to-end system that (1) captures \emph{process, file, network and privilege}
behaviour of Kubernetes pods with eBPF, (2) attaches \emph{Pod, container, namespace
and workload} identity to every event, (3) combines a supervised multi-class model
with interpretable rules to produce a \emph{risk score, severity and MITRE technique},
and (4) is evaluated honestly on synthetic, real-world and live-cluster data including
latency and resource overhead. This project addresses that gap.
